\section{Clone a Graph}
\label{sec:graph-clone-graph}

\problemstatement{Given a reference to a node in a connected
undirected graph (where each node stores a value and a list of
neighbor pointers), return a \textbf{deep copy} (clone) of the entire
graph -- a structurally identical graph made of entirely new node
objects.}

\begin{patternbox}
\emph{``Copy every node and every connection, visiting each node
exactly once''} is a traversal (DFS or BFS) augmented with a
\textbf{hash map from original node to its clone} -- the map serves
double duty: it remembers which nodes have already been cloned (acting
as the traversal's \textit{visited} check) and it provides the correct
clone object to link into a neighbor's adjacency list.
\end{patternbox}

\subsection*{Why a hash map, not a simple \textit{visited} array}

Every traversal in this chapter needs to track which nodes have been
visited -- but previous problems used array indices (city numbers,
grid coordinates) as keys, which a simple boolean array handles
directly. Here, nodes are objects with no natural array index, and
\emph{visiting} a node is inseparable from \emph{creating its clone} --
so a hash map from original node pointer to cloned node pointer plays
both roles: checking $\textit{map.find(node)} \neq \textit{map.end()}$
is the visited check, and $\textit{map}[\textit{node}]$ retrieves the
already-built clone to attach as a neighbor.

\subsection*{Algorithm (DFS)}

\begin{enumerate}
  \item Define $\textit{cloneNode}(\textit{node})$: if $\textit{node}$
        is already in the map, return its clone immediately (this
        breaks cycles -- without this check, a cyclic graph would
        recurse forever).
  \item Otherwise, create a new node with the same value, store it in
        the map \textbf{before} recursing into neighbors (exactly like
        marking visited before recursing in plain DFS), then for each
        original neighbor, recursively clone it and append the result
        to the new node's neighbor list.
  \item Return $\textit{cloneNode}(\textit{startNode})$.
\end{enumerate}

\subsection*{C++ implementation}

\begin{lstlisting}
#include <bits/stdc++.h>
using namespace std;

struct Node {
    int val;
    vector<Node*> neighbors;
    Node(int v) : val(v) {}
};

Node* cloneNode(Node* node, unordered_map<Node*, Node*>& cloneMap) {
    if (cloneMap.find(node) != cloneMap.end()) {
        return cloneMap[node];
    }

    Node* clone = new Node(node->val);
    cloneMap[node] = clone;

    for (Node* neighbor : node->neighbors) {
        clone->neighbors.push_back(cloneNode(neighbor, cloneMap));
    }
    return clone;
}

Node* cloneGraph(Node* startNode) {
    if (startNode == nullptr) return nullptr;
    unordered_map<Node*, Node*> cloneMap;
    return cloneNode(startNode, cloneMap);
}

int main() {
    Node* a = new Node(1);
    Node* b = new Node(2);
    Node* c = new Node(3);
    a->neighbors = {b, c};
    b->neighbors = {a, c};
    c->neighbors = {a, b};

    Node* clonedA = cloneGraph(a);
    cout << clonedA->val << " has " << clonedA->neighbors.size()
         << " neighbors" << endl; // 1 has 2 neighbors
    cout << (clonedA == a ? "same object" : "different object") << endl; // different object
    return 0;
}
\end{lstlisting}

\subsection*{Dry run}

\dryrun{Take a triangle graph: nodes $1,2,3$, each connected to the
other two. Start cloning from node $1$.}

$\textit{cloneNode}(1)$: not in map, create clone $1'$, store
$\textit{map}[1]=1'$. Recurse into neighbor $2$:
$\textit{cloneNode}(2)$: not in map, create clone $2'$, store
$\textit{map}[2]=2'$. Recurse into $2$'s neighbors, $1$ and $3$.
$\textit{cloneNode}(1)$: \textbf{already in map}, return $1'$
immediately (no infinite recursion, even though $1$ and $2$ point back
to each other). $\textit{cloneNode}(3)$: not in map, create clone
$3'$, store $\textit{map}[3]=3'$; recurse into $3$'s neighbors $1,2$,
both already in map, return $1',2'$ immediately. Back at $2'$:
neighbors $=[1',3']$. Back at $1'$: continue to neighbor $3$, already
in map, return $3'$. Final: $1'$'s neighbors $=[2',3']$ -- a complete,
independent clone of the original triangle.

\begin{complexitybox}
\textbf{Time Complexity.} $O(V+E)$ -- every node is cloned exactly
once (guarded by the map check), and every edge is examined exactly
once (from each direction).
\textbf{Space Complexity.} $O(V)$ for the hash map and the recursion
stack.
\end{complexitybox}

\begin{takeawaybox}
Whenever a traversal needs to \textbf{build something new} while
visiting (a clone, a transformed copy, a mapping to a different
representation) rather than simply recording or counting, a hash map
keyed by the original node typically replaces the plain
\textit{visited} array -- the map's presence-check still prevents
revisiting and infinite loops, while also handing back the
already-built result the moment it's needed again.
\end{takeawaybox}
